\documentclass[12pt,a4paper]{article}
\usepackage[margin=2.4cm]{geometry}
\usepackage{amsmath,amssymb,bm,mathtools}
\usepackage{fontspec}
\setmainfont{TeX Gyre Termes}
\setsansfont{TeX Gyre Heros}
\setmonofont{Latin Modern Mono}[Scale=0.9]
\usepackage{graphicx,float,caption,subcaption,booktabs,longtable,array,multirow,tabularx}
\usepackage{xcolor}
\usepackage{setspace}
\usepackage{lineno}
\usepackage{microtype}
\usepackage{xurl}
\usepackage{siunitx}
\usepackage{enumitem}
\usepackage{fancyhdr}
\usepackage{titlesec}
\titleformat{\paragraph}[runin]{\normalfont\itshape}{}{0pt}{}
\titlespacing*{\paragraph}{0pt}{1.1ex plus .2ex}{0.7em}
\usepackage[authoryear,round,semicolon,sort&compress]{natbib}
\definecolor{lk}{RGB}{0,70,140}
\usepackage[colorlinks=true,linkcolor=lk,citecolor=lk,urlcolor=lk,breaklinks=true,pdfborder={0 0 0}]{hyperref}
\usepackage[nameinlink,capitalise]{cleveref}
\sisetup{detect-all,separate-uncertainty=false}
\setstretch{1.5}
\setlength{\parskip}{3pt}
\captionsetup{font=small,labelfont=bf,justification=centering,singlelinecheck=true}
\titleformat{\section}{\large\bfseries}{\thesection.}{0.5em}{}
\titleformat{\subsection}{\normalsize\bfseries}{\thesubsection}{0.5em}{}
\titleformat{\subsubsection}{\normalsize\itshape}{\thesubsubsection}{0.5em}{}
\pagestyle{fancy}\fancyhf{}\fancyfoot[C]{\small\thepage}\renewcommand{\headrulewidth}{0pt}
\setlength{\headheight}{14pt}
\bibliographystyle{plainnathl}
\newcommand{\doihl}[1]{doi: \href{https://doi.org/#1}{\nolinkurl{#1}}}
\newcommand{\dd}{\mathrm{d}}
\newcommand{\ET}{\mathit{ET}}
\newcommand{\Ent}{\mathit{ET}_{\mathrm{nt}}}
\newcommand{\Ks}{K_{\mathrm{s}}}
\newcommand{\Se}{S_{\mathrm{e}}}
\newcommand{\TAW}{\mathit{TAW}}
\newcommand{\repourl}{https://github.com/halilunaziru73-creator/Coupled-Water-Fluxes-Pore-Scale-Dynamics-Precision-Irrigation-Mulching-Feedback-Loops}
\newcommand{\repo}{\href{\repourl}{repository}}
\newcommand{\datalink}{\href{\repourl}{Data}}
\newcommand{\codelink}{\href{\repourl/tree/main/Manuscripts}{Code}}
\newcommand{\softlink}{\href{\repourl/tree/main/Software/SWAFL_Desktop}{Software}}
\newcommand{\manuallink}{\href{\manualurl}{Manual}}
\newcommand{\downloadlink}{\href{\exeurl}{SWAFL.exe}}
\newcommand{\figdir}{../figures}
\newcommand{\manualurl}{https://github.com/halilunaziru73-creator/Coupled-Water-Fluxes-Pore-Scale-Dynamics-Precision-Irrigation-Mulching-Feedback-Loops/blob/main/Software/SWAFL_Desktop/docs/SWAFL_Manual.pdf}
\newcommand{\exeurl}{https://github.com/halilunaziru73-creator/Coupled-Water-Fluxes-Pore-Scale-Dynamics-Precision-Irrigation-Mulching-Feedback-Loops/releases/latest/download/SWAFL.exe}
\newcommand{\authorblock}{%
\begin{center}
{\large\textbf{Naziru Halilu}}$^{1,2,*}$\\[3pt]
{\small $^{1}$ Agricultural University of Athens, Athens, Greece\\
$^{2}$ Public University of Navarre, Pamplona, Spain\\
ORCID: \href{https://orcid.org/0009-0006-2324-9599}{0009-0006-2324-9599}\\
$^{*}$ Correspondence: \href{mailto:halilunaziru73@gmail.com}{halilunaziru73@gmail.com}}
\end{center}}
\newcommand{\codeavail}{%
\datalink{} (the complete Daisy ensemble outputs), \codelink{} (analysis and figure scripts) and \softlink{} (the SWAFL Desktop source) are in the project repository; the \manuallink{} (93 pages) can be reviewed or downloaded as a PDF. The dataset (all 10\,190 runs, raw and derived) is archived on Zenodo, \doihl{10.5281/zenodo.23043626}, and the ensemble table used here is also stored in the repository as \texttt{Manuscripts/reproducibility/data/daisy\_ensemble.csv.gz}. The folder \texttt{Manuscripts/reproducibility} is a self-contained package: pinned software versions (Python 3.13.16, NumPy 2.5.3, SciPy 1.18.1, pandas 3.0.5, scikit-learn 1.9.1, Matplotlib 3.11.2, Node.js 22.22.0), fixed random seeds in every stochastic step, SHA-256 checksums, and one command (\texttt{run\_all.sh}) that regenerates all results and figures from the ensemble and then compares every regenerated value with the archived result and with the headline numbers printed in the manuscripts. Each table and figure names the script that produced it (Appendix~\ref{app:repro}). Daisy itself was not re-run for this package, no field measurements are used, and three small intermediate arrays and the interface screenshots are archived rather than regenerated; these limits are listed in \texttt{README\_REPRODUCE.md}.}

\usepackage{pdflscape}
\setlength{\emergencystretch}{3em}

\hypersetup{pdftitle={A Verified Bridge between a Process-Based Soil-Crop Model and Field Decision Support},pdfauthor={Naziru Halilu}}
\begin{document}
\linenumbers
\begin{titlepage}
\vspace*{-1.2cm}
\begin{center}\begin{minipage}{0.92\textwidth}\centering{\Large\bfseries From Simulation to Decision: A Verified Bridge between the Daisy Model and Field Decision-Support Layers\par}\end{minipage}\end{center}
\vspace{0.4cm}
\authorblock
\vspace{0.4cm}
{\small\setlength{\parindent}{0pt}\setlength{\parskip}{0pt}
\begin{center}\textbf{\normalsize Abstract}\end{center}\vspace{-0.4em}
\noindent Process-based models answer the questions of agronomy and irrigation only when their output reaches the person who has to decide, in a form that person can use. Decision tools, in turn, are only as trustworthy as the simulations behind them. This article describes and tests a bridge that connects the Daisy soil--plant--atmosphere model to the decision layers of an offline suite (SWAFL Desktop 1.1.0) in three ways: through a pre-computed ensemble of 10\,190 Daisy runs that a fast meta-model (Engine 1) reproduces; through the live runs that a user makes in a graphical Daisy workbench, whose logs are parsed, summarised by season and handed to two decision layers (a decision table and the fate of a chemical) and matched to the other layers that answer the same question; and through the weather of the project, which is passed back to Daisy. The bridge has five parts (a log parser, a seasonal summariser, a soil map, a question-to-layer map and a shared live object) and one rule: a decision layer draws a Daisy result only from a Daisy run, and otherwise says that there is none. Each part was tested against the files of the Daisy project. The parsers read the six bundled logs (11\,665 rows in the largest) with a relative difference of $3\times10^{-15}$ from an independent reader; the seasonal nitrogen summaries reproduce the log columns to the 0.1\,kg\,ha$^{-1}$ rounding of the summary; the chemical budget of an annual tracer log closes to 0.12\,kg\,ha$^{-1}$ per season (0.02\,\% of the 2\,290\,kg\,ha$^{-1}$ applied); and the denitrification function that the bridge shares with Daisy reproduces 792 stored Daisy outputs to $3\times10^{-15}$ after a defect, found by this test, was corrected. The three routes differ by seven orders of magnitude in time per scenario (a median of 145\,s per Daisy run, 17.5 days for the ensemble, about $2\times10^{-5}$\,s per scenario for Engine 1) and in what they can answer: the ensemble and Engine 1 give fast and ranked answers for the soils and weather of the study, the live route gives the answer for the user's own soil and weather. A complete end-to-end run, with Daisy executing a set-up and the layers filling, could not be made in the environment of this study, and no field data are used; both limits are stated, together with the unmapped soils of the bridge.

\medskip\noindent \textbf{Keywords.} model coupling; decision support; Daisy; meta-model; nitrogen budget; verification; irrigation scheduling; reproducible research\par}
\end{titlepage}

\section{Introduction}
A simulation model is a function from a description of a field to a set of time series. A decision is a choice between management options for a particular field. The distance between the two is crossed by operations that are rarely described or tested: the reading of the output, its reduction to a few numbers per season, the matching of the simulated soil to the soil of the field, the choice of the output that answers the question, and the presentation of the answer next to its alternatives. When a decision tool is built on a process model, these operations are the tool.

Two strategies are common. In the first, the model is run offline over a design of scenarios and a fast statistical model (a meta-model, emulator or surrogate) is fitted to the results \citep{Sacks1989,Kennedy2001,Razavi2012}; the decision tool then evaluates the emulator. This is fast and needs no model installation, but it can only answer for the conditions of the design. In the second, the tool runs the model on demand for the user's own conditions, which is general but slow and needs the model. Both have a place in a tool for irrigation planning, and the tool described here uses both: a meta-model of a published ensemble of 10\,190 Daisy runs \citep{HaliluData} is Engine~1 of the suite (Paper~3), and a graphical workbench for Daisy (Paper~4) runs Daisy for the user and hands the results on.

The question of this article is how the hand-over should be built so that it can be checked. We describe the components of the bridge, give the rule by which the decision layers use them, and test each component against files that the Daisy project itself provides. The tests are a part of the result because a bridge that mistranslates a log is worse than no bridge: it produces numbers that look like Daisy output.

\paragraph{Contributions.} (1) A bridge with five small parts and a single rule that keeps the origin of every Daisy number visible. (2) A set of tests of each part against the files of the Daisy project (logs of nitrogen, a tracer and water, and 792 stored outputs of the denitrification function). (3) A comparison of the three routes from Daisy to a decision (ensemble, meta-model, live run) in time and in what they can answer. (4) A defect of the shared denitrification function that the tests found, and its correction.

\section{The bridge}
\subsection{Parts}
\Cref{fig:bridge} shows the bridge. A scenario or an experiment of the workbench is turned into a Daisy set-up and run by Daisy 7.1.14. Daisy writes its logs as tab-separated files (\texttt{.dlf}) with a header, a row of names, a row of units and one row per output interval. The bridge has five parts.

\begin{description}[leftmargin=1.5em,style=nextline]
\item[Log parser (\texttt{DaisyIO}, \texttt{DaisySum.parseDlf}).] Reads a log into a list of names and rows. Numeric columns become numbers; the crop name stays text. The first row of an annual log is written at the start of the simulation and is not a season; the summariser drops it.
\item[Seasonal summariser (\texttt{DaisySum.summarise}).] Reduces the logs of a run to one record with a value per season: harvest of the crop (dry matter of the storage organ, water stress, nitrogen stress), precipitation, irrigation, actual and potential evapotranspiration, runoff, percolation below 1\,m (matrix and biopore flow), and the nitrogen terms of \cref{eq:nsum}. For a run of a chemical it adds the budget of \cref{eq:chem}.
\item[Soil map (\texttt{SOILMAP}).] Maps the soil of a Daisy run to the soil of the suite (for example \texttt{usda:siltloam} to Silt Loam, \texttt{S07}), so that the other layers show the same soil. Eight soils are mapped (sand, loamy sand, sandy loam, loam, silt loam, sandy clay loam, clay loam, silty clay loam); the two coarse soils that are special to the study (Coarse Loam, Coarse Sand) and an own soil table have no counterpart, and the layers then keep the soil of the project.
\item[Question map (\texttt{MATCH}).] For each of five questions (irrigation, nitrogen, climate, a chemical, the soil) lists the layers that answer it with a sentence on what each shows (\cref{tab:match}). The hand-over tab of the workbench opens the first of them.
\item[Live object (\texttt{DAISYLIVE}).] Holds, in the engine of the suite, the label of the weather and soil of the run, the summarised runs, their settings (irrigation, nitrogen factor, rain factor, temperature shift) and the mapped soil. The decision layers read only this object.
\end{description}

\begin{figure}[H]\centering
\includegraphics[width=\textwidth]{\figdir/p5/fig01_bridge.pdf}
\caption{The bridge from a Daisy run to a decision. Solid arrows are the path of a run; the dashed arrow at the left is the return path of the weather of the project to Daisy.}\label{fig:bridge}
\end{figure}

\begin{table}[H]\centering\small
\caption{Questions of the workbench and the layers that answer them.}\label{tab:match}
\begin{tabularx}{\textwidth}{>{\raggedright\arraybackslash}p{3.2cm}X}
\toprule Question & Layers\\\midrule
Irrigation: does it pay? & Daisy decision table (\texttt{dsum}); reliability over 20 seasons (\texttt{yby}); full water balance (\texttt{wb2}); sensor and trigger (\texttt{v1})\\
Nitrogen: how much fertiliser? & decision table; nitrogen budget (\texttt{nb})\\
Climate: warmer, drier, wetter & decision table; reliability over 20 seasons\\
A chemical: sorption and decay & fate of a chemical (\texttt{chem})\\
The soil: how much does texture matter? & soil hydraulic drivers (\texttt{v7}); profile evolution (\texttt{v2}); decision table\\
Own set-up & decision table (own runs next to the others when the names match)\\\bottomrule
\end{tabularx}
\end{table}

\subsection{What the summariser computes}
For season $s$ the nitrogen terms are sums of columns of the annual nitrogen log (kg\,N\,ha$^{-1}$):
\begin{equation}\label{eq:nsum}
\begin{aligned}
N^{\mathrm{fert}}_s&=N_{\mathrm{surf}}+N_{\mathrm{soil}}+N_{\mathrm{org}}, & N^{\mathrm{up}}_s&=N_{\mathrm{uptake}},\\
N^{\mathrm{den}}_s&=N_{\mathrm{denit}}+N_{\mathrm{N_2O\,nitrif}}, & N^{L}_s&=N_{\mathrm{matrix\,leach}}+N_{\mathrm{biopore\,leach}}\ \ (\text{at }100\ \mathrm{cm}),
\end{aligned}
\end{equation}
together with the mineralisation, the deposition and the volatilisation. The chemical budget uses the annual tracer log (kg\,ha$^{-1}$; the log is in g):
\begin{equation}\label{eq:chem}
\Delta M_s = A_s - U_s - L_s - R_s - X_s,\qquad L_s=L_{\mathrm{matrix}}+L_{\mathrm{biopore}}+L_{\mathrm{soil\,drain}}+L_{\mathrm{biopore\,drain}},
\end{equation}
where $A$ is the spray, $U$ the crop uptake, $R$ the runoff, $X$ the decay (soil, surface and litter), and $\Delta M$ the change of the amount in the soil from one season to the next. A depth profile of the concentration is added from the daily concentration log. The weather and soil of the run are described by a one-line label, shown at the top of every Daisy layer (``Daisy 7.1.14, run in the Daisy workbench: weather Taastrup, Denmark; soil loam; years 1979 to 2000'').

\subsection{The rule of origin}
A decision layer that shows a Daisy quantity draws it from a live run or not at all. Before a run, the two layers of Daisy output (the decision table and the fate of a chemical) show the sentence ``no Daisy run yet'', say what is needed to fill them, and offer a button to the scenarios of the workbench; they do not substitute bundled or estimated numbers. The other layers of the suite, which are driven by Engine~1 and Engine~2, say in their own header that they show the meta-model and the daily model. The rule is a small feature, but it settles a question that every coupled tool has to answer: when a number is on the screen, is it from the model it is labelled with? In a test with the application opened afresh (\texttt{p5\_empty\_state.js}) both layers showed this sentence and no figure, and the engine held no live object.

\section{Methods of the tests}
\paragraph{Parsers.} The six Daisy logs bundled with the suite (annual field nitrogen, two runs; annual tracer; daily water flux; the nitrogen log of a calibration example; a soil chemical log with 11\,665 rows) were read by the parser of the bridge and, independently, by \texttt{pandas}. For every numeric column the sum was compared; the largest relative difference is reported.

\paragraph{Nitrogen summaries.} The two annual nitrogen logs hold ten seasons of a rotation, with the same fertiliser and different set-ups (runs 3b and 4b). For each season the summariser was compared with the sum of the corresponding log columns of \cref{eq:nsum}, and the closure column of the log (Daisy's own error term) was read.

\paragraph{Chemical budget.} The annual tracer log was passed through the summariser with chemistry, and the residual of \cref{eq:chem} was computed for each of the ten seasons from the summarised amounts.

\paragraph{Denitrification.} The function that the workbench shares with Daisy was compared with the 792 inputs and stored outputs of the Daisy denitrification test (Paper~4 describes the function and its defect).

\paragraph{Time and agreement of the three routes.} The Daisy time of the ensemble and of its experiments was summed from the recorded run times (\texttt{run\_seconds}); the time of Engine~1 and its agreement with Daisy on held-out runs are those of Paper~3.

\section{Results}
\subsection{Parsing and summarising}
The six logs were read with relative differences between the two readers of at most $3\times10^{-15}$ (largest log) and 0 for the calibration log, which is a single row. The seasonal nitrogen summaries differ from the log columns by at most 0.049\,kg\,ha$^{-1}$ in the uptake, 0.049 in the mineralisation, 0.005 in the denitrification and the leaching, and 0 in the fertiliser; these differences are the rounding of the summary to 0.1 (nitrogen) or 0.01\,kg\,ha$^{-1}$, and the Daisy error column is 0 in all 20 seasons. In the chemical budget the residual of \cref{eq:chem} is 0.12\,kg\,ha$^{-1}$ or less in each of ten seasons and 0.47 over the ten, 0.02\,\% of the 2\,290\,kg\,ha$^{-1}$ applied, which are the rounding of the log values. The terms of the budget can be trusted to the precision of the log.

The two nitrogen runs show what the layers are for. Their mean fertiliser is the same (229\,kg\,N\,ha$^{-1}$ per season), the crop takes up 341 and 339\,kg\,N\,ha$^{-1}$, but the leaching below 1\,m is 4.4 in run 3b and 17.0 in run 4b, a factor of 3.9, and the denitrification is 15.8 and 14.5 (\cref{tab:nrun}). A person who compared only the uptake would see no difference. The bridge reports all terms of the budget side by side, and the nitrogen layer shows them season by season.

\begin{table}[H]\centering\small
\caption{Mean seasonal nitrogen budget of the two bundled annual nitrogen logs as reported by the bridge (kg\,N\,ha$^{-1}$, ten seasons).}\label{tab:nrun}
\begin{tabular}{lrrrrrr}\toprule
Run & Fertiliser & Mineralisation & Deposition & Crop uptake & Denitrification & Leaching\\\midrule
3b & 229.0 & 136.0 & 13.2 & 341.5 & 15.8 & 4.4\\
4b & 229.0 & 134.9 & 13.2 & 338.5 & 14.5 & 17.0\\\bottomrule
\end{tabular}
\end{table}

\subsection{The denitrification defect}
The first comparison of the shared denitrification function with the 792 stored outputs gave a median relative error of 40\,\%, because the rescaling of the carbon dioxide production for the pressure head was applied to inputs that Daisy had already rescaled. After the correction the largest relative difference is $2.9\times10^{-15}$ (N$_2$) and $1.7\times10^{-15}$ (N$_2$O) (Paper~4). The function is used in the workbench to draw the response curves and, through the bridge, in the reading of the gaseous nitrogen losses.

\subsection{Three routes from Daisy to a decision}
The study ran 10\,190 simulations in 17 experiments in 17.5 days of Daisy time (mean 148\,s, median 145\,s per run). Evaluating Engine~1 for $10^5$ scenarios takes about 2\,s (Paper~3), that is $2\times10^{-5}$\,s per scenario, a factor of $7\times10^{6}$ against the median Daisy run. The compact experiments of the workbench stand in between, 368 runs in about 15 hours (Paper~4). The three routes are not alternatives of one kind but answer different questions (\cref{tab:routes}, \cref{fig:econ}).

\begin{table}[H]\centering\small
\caption{The three routes from Daisy to a decision.}\label{tab:routes}
\begin{tabularx}{\textwidth}{>{\raggedright\arraybackslash}p{3.0cm}XXX}
\toprule & Ensemble and Engine 1 & Compact experiment (workbench) & Own live run\\\midrule
Soil and weather & ten soils, one weather record, the design of the study & as chosen in the workbench (library or own soil, bundled or own weather) & own soil table and weather\\
Time & microseconds per scenario & about 0.2 to 4 hours per experiment & minutes per run\\
Needs Daisy & no & yes (Windows) & yes (Windows)\\
Answers & ranking of options, response curves, regime maps within the design & the shape of a response for the chosen conditions & the budget of one set-up, for the layers of the suite\\
Main limit & fixed design; offset of absolute levels (Paper~3) & coarse levels (Paper~4) & one run, one set-up\\\bottomrule
\end{tabularx}
\end{table}

\begin{figure}[H]\centering
\includegraphics[width=\textwidth]{\figdir/p5/fig04_economics.pdf}
\caption{(a) Daisy wall-clock time of the 17 experiments of the study (total 17.5 days). (b) Time of the study, of the compact designs of the workbench and of Engine 1 for $10^5$ scenarios (logarithmic scale).}\label{fig:econ}
\end{figure}

\begin{figure}[H]\centering
\includegraphics[width=\textwidth]{\figdir/p5/fig02_flowscreens.jpg}
\caption{The hand-over in the application. (a) The last step of the guided path of the workbench: five questions, each with the layers that answer it. (b) The Daisy decision table with the labels of the run (rainfed and irrigated, half and 30\,\% more fertiliser, warmer and drier climates). (c) The fate of a chemical for a mobile, a sorbing and a sorbing and degrading tracer.}\label{fig:flow}
\end{figure}

\subsection{Agreement of the fast route with Daisy}
The agreement of Engine~1 with Daisy was tested on 1\,000 runs of the sensor, trigger and rain experiment (E12, coarse loam) that were not used to build its lookups: the ranking of the trigger, sensor-depth and rainfall options is preserved (Spearman 0.64 to 0.97) and the absolute level has a diagnosed offset (Paper~3). This is the evidence that the first route can be used to choose between options. It is not evidence that the second and third routes agree with the first for a soil outside the study; for those the user is told, in the header of the layer, which weather and soil were used.

\section{Discussion}
\paragraph{Why the rule of origin matters.} A bridge between a model and a tool has two kinds of risk. One is translation (a column read wrongly, a unit lost); the tests above address it. The other is substitution: a layer that shows a number from a different source under the name of the model. The rule that a Daisy layer shows a Daisy number or nothing removes the second kind by design, and the label of the run keeps it visible.

\paragraph{What the tests did not cover.} (i) Daisy was not run in this study, so the whole chain from a set-up to a filled layer was tested in parts and not as one. The parts are the places where a defect can hide, and the one defect found (the denitrification function) was in a part. An end-to-end test with a Daisy installation is the first thing to do. (ii) The bundled logs have ten seasons of one rotation; the logs of a year with a different crop, a long fallow or irrigation have columns that were not exercised. (iii) The two coarse soils of the study and own soil tables are not mapped, so the layers of the suite keep the soil of the project in those cases. (iv) The agreement of the first route with Daisy was tested against Daisy runs of the study itself, not against measurements. (v) The summariser reduces a run to seasons; a question that depends on the day (the date of the first irrigation) needs the daily logs, which the workbench plots but the layers do not read.

\paragraph{Generalisation.} The design is not specific to Daisy: a parser of the logs of another model, a summariser to seasonal terms, a map of the soils, a map of questions to layers and one shared object are the parts of any such bridge. The rule of origin and the tests against the files of the model's own project are what make it checkable.

\section{Conclusions}
A bridge from a process model to decision layers can be small, rule-bound and tested part by part. The bridge from Daisy to the SWAFL layers reads the logs exactly, summarises them to the precision of the log, closes the chemical budget to 0.02\,\% of the amount applied, and shares with Daisy a denitrification function that reproduces 792 stored outputs after the correction of one defect. It offers three routes that differ by seven orders of magnitude in time. The next steps are an end-to-end test on a machine with Daisy and a comparison of the live layers with measured seasons.

\section*{Software, data and code availability}
\suiteavail
\section*{Author contributions}
N.H. designed the bridge, wrote the software and the tests, and wrote the article.
\section*{Funding}
No external funding was received for this work. [To be completed by the author if applicable.]
\section*{Competing interests}
The author declares no competing interests.

\bibliography{refs}

\appendix
\section{Checks and scripts}\label{app:repro5}
\begin{table}[H]\centering\small
\caption{Check to script map (scripts in \texttt{Manuscripts/reproducibility/p456}).}
\begin{tabularx}{\textwidth}{>{\raggedright\arraybackslash}p{4.6cm}>{\raggedright\arraybackslash}p{4.0cm}X}
\toprule Check & What is compared & Script $\to$ output\\\midrule
Parsers and nitrogen summaries & bridge against \texttt{pandas} & \texttt{p5\_bridge\_checks.js}, \texttt{p5\_bridge\_checks.py} $\to$ \texttt{out\_p5\_bridge.json}\\
Chemical budget & residual of \cref{eq:chem} & \texttt{p5\_chem\_closure.js} $\to$ \texttt{out\_p5\_chem.json}\\
Denitrification & 792 stored outputs & \texttt{p5\_bridge\_checks.js} $\to$ \texttt{out\_p5\_bridge\_js.json}\\
Run time & \texttt{run\_seconds} of the ensemble & \texttt{p4\_design\_fidelity.py} $\to$ \texttt{out\_p4\_runtime.csv}\\
Empty state of the Daisy layers & application opened afresh & \texttt{p5\_empty\_state.js}\\
Figures & & \texttt{p456\_figs\_schematics.py}, \texttt{p456\_figs\_data.py}\\\bottomrule
\end{tabularx}
\end{table}
\end{document}
